% !TEX root = ../mainsempro.tex
% Salinan kerja Bab II khusus revisi mainsempro; sumber awal: chapters/tinjauan.tex

\chapter{Tinjauan Pustaka}
\label{chap:tinjauan}

% ============================================================
\section{Geometri Ruang-Waktu dan Tensor Metrik}
\label{sec:geometri-metrik}

Dalam relativitas umum, gravitasi direpresentasikan melalui geometri
ruang-waktu empat dimensi. Struktur geometris tersebut dinyatakan oleh
tensor metrik $g_{\mu\nu}$, yang menentukan interval antara dua peristiwa
yang berdekatan \parencite{Schutz2009FirstCourse,Carroll2004SpacetimeGeometry}:
\begin{equation}
    ds^2 = g_{\mu\nu}\,dx^\mu dx^\nu.
    \label{eq:interval-umum}
\end{equation}

Dalam penelitian ini digunakan konvensi tanda $(-,+,+,+)$. Pada
ruang-waktu datar, tensor metrik direduksi menjadi metrik Minkowski,
\begin{equation}
    \eta_{\mu\nu}
    =
    \operatorname{diag}(-1,1,1,1),
    \label{eq:minkowski-metric}
\end{equation}
sehingga interval ruang-waktu dapat dituliskan sebagai
\begin{equation}
    ds^2
    =
    -c^2dt^2
    +dx^2+dy^2+dz^2.
    \label{eq:minkowski-interval}
\end{equation}

Berbeda dengan ruang-waktu datar, pada ruang-waktu melengkung komponen
metrik secara umum bergantung pada koordinat,
\begin{equation}
    g_{\mu\nu}=g_{\mu\nu}(x^\alpha).
\end{equation}
Ketergantungan tersebut menyebabkan bentuk geometri berubah dari satu
titik ke titik lainnya.

Selain menentukan interval, tensor metrik digunakan untuk menghubungkan
komponen kontravarian dan kovarian suatu vektor. Untuk vektor
$V^\mu$, penurunan indeks dilakukan melalui
\begin{equation}
    V_\mu = g_{\mu\nu}V^\nu,
    \label{eq:lowering-index}
\end{equation}
sedangkan operasi sebaliknya menggunakan metrik invers $g^{\mu\nu}$,
\begin{equation}
    V^\mu = g^{\mu\nu}V_\nu.
    \label{eq:raising-index}
\end{equation}
Metrik invers didefinisikan melalui
\begin{equation}
    g^{\mu\alpha}g_{\alpha\nu}
    =
    \delta^\mu_{\ \nu},
    \label{eq:inverse-metric-def}
\end{equation}
sehingga secara matriks berlaku
\begin{equation}
    [g^{\mu\nu}]
    =
    [g_{\mu\nu}]^{-1}.
\end{equation}

Karena komponen metrik pada umumnya bergantung pada posisi, turunan
$\partial_\lambda g_{\mu\nu}$ tidak selalu bernilai nol. Akibatnya,
turunan biasa terhadap komponen vektor belum cukup untuk
mendeskripsikan perubahan suatu vektor di sepanjang ruang-waktu.

Suatu vektor dapat dituliskan sebagai
\begin{equation}
    \mathbf{V}=V^\mu\mathbf{e}_\mu,
\end{equation}
dengan $\mathbf{e}_\mu$ menyatakan vektor basis. Turunan terhadap
koordinat $x^\nu$ menghasilkan
\begin{equation}
    \partial_\nu\mathbf{V}
    =
    (\partial_\nu V^\mu)\mathbf{e}_\mu
    +
    V^\mu(\partial_\nu\mathbf{e}_\mu).
    \label{eq:vector-derivative-basis}
\end{equation}
Persamaan~\eqref{eq:vector-derivative-basis} menunjukkan bahwa perubahan
vektor terdiri atas dua bagian, yaitu perubahan komponen $V^\mu$ dan
perubahan vektor basis $\mathbf{e}_\mu$.

Perubahan basis dapat dinyatakan menggunakan koefisien koneksi,
\begin{equation}
    \partial_\nu\mathbf{e}_\mu
    =
    \Gamma^\lambda_{\mu\nu}\mathbf{e}_\lambda,
    \label{eq:basis-connection}
\end{equation}
sehingga turunan vektor menjadi
\begin{equation}
    \partial_\nu\mathbf{V}
    =
    \left(
        \partial_\nu V^\lambda
        +
        \Gamma^\lambda_{\mu\nu}V^\mu
    \right)\mathbf{e}_\lambda.
\end{equation}
Kombinasi di dalam tanda kurung didefinisikan sebagai turunan
kovarian vektor kontravarian,
\begin{equation}
    \nabla_\nu V^\lambda
    =
    \partial_\nu V^\lambda
    +
    \Gamma^\lambda_{\mu\nu}V^\mu.
    \label{eq:covariant-derivative}
\end{equation}

Dalam relativitas umum digunakan koneksi Levi-Civita, yaitu koneksi
yang kompatibel dengan metrik,
\begin{equation}
    \nabla_\lambda g_{\mu\nu}=0,
    \label{eq:metric-compatibility}
\end{equation}
dan bebas torsi,
\begin{equation}
    \Gamma^\rho_{\mu\nu}
    =
    \Gamma^\rho_{\nu\mu}.
    \label{eq:torsion-free}
\end{equation}
Dari kedua kondisi tersebut diperoleh simbol Christoffel jenis kedua
\parencite{Schutz2009FirstCourse,Carroll2004SpacetimeGeometry}:
\begin{equation}
    \Gamma^\rho_{\mu\nu}
    =
    \frac{1}{2}g^{\rho\sigma}
    \left(
        \partial_\mu g_{\sigma\nu}
        +
        \partial_\nu g_{\sigma\mu}
        -
        \partial_\sigma g_{\mu\nu}
    \right).
    \label{eq:christoffel-def}
\end{equation}

Persamaan~\eqref{eq:christoffel-def} menunjukkan bahwa simbol
Christoffel ditentukan oleh tensor metrik, metrik invers, dan turunan
pertama komponen metrik. Dengan demikian, hubungan matematis antara
metrik dan koneksi dapat diringkas sebagai
\begin{equation}
    g_{\mu\nu}
    \longrightarrow
    g^{\mu\nu}
    \longrightarrow
    \partial_\alpha g_{\mu\nu}
    \longrightarrow
    \Gamma^\rho_{\mu\nu}.
\end{equation}

Secara fisis, tensor metrik menentukan struktur pengukuran ruang dan
waktu, sedangkan simbol Christoffel digunakan untuk membandingkan
vektor pada titik-titik yang berbeda. Namun, simbol Christoffel bukan
ukuran intrinsik kelengkungan karena nilainya dapat tidak nol akibat
pemilihan sistem koordinat. Kelengkungan intrinsik ruang-waktu baru
dinyatakan melalui tensor Riemann, yang dibahas pada bagian berikutnya
\parencite{Schutz2009FirstCourse,Carroll2004SpacetimeGeometry}.


% ============================================================
\section{Kelengkungan Ruang-Waktu dan Persamaan Medan Einstein}
\label{sec:kelengkungan-efe}

Kelengkungan intrinsik ruang-waktu dinyatakan oleh tensor Riemann. Salah satu cara untuk memperkenalkan tensor ini adalah melalui komutator dua turunan kovarian yang bekerja pada vektor $V^\rho$ \parencite{Schutz2009FirstCourse,Carroll2004SpacetimeGeometry}:
\begin{equation}
    \left(\nabla_\mu\nabla_\nu-\nabla_\nu\nabla_\mu\right)V^\rho
    =R^\rho_{\ \sigma\mu\nu}V^\sigma.
    \label{eq:riemann-commutator}
\end{equation}
Dengan konvensi indeks yang digunakan dalam penelitian ini, tensor Riemann dituliskan sebagai
\begin{equation}
    R^\rho_{\ \sigma\mu\nu}
    =\partial_\mu\Gamma^\rho_{\nu\sigma}
    -\partial_\nu\Gamma^\rho_{\mu\sigma}
    +\Gamma^\rho_{\mu\lambda}\Gamma^\lambda_{\nu\sigma}
    -\Gamma^\rho_{\nu\lambda}\Gamma^\lambda_{\mu\sigma}.
    \label{eq:riemann}
\end{equation}
Karena simbol Christoffel dibangun dari turunan pertama tensor metrik, struktur tensor Riemann secara skematis dapat dinyatakan sebagai
\begin{equation}
    R^\rho_{\ \sigma\mu\nu}
    \sim \partial\Gamma+\Gamma\Gamma
    \sim \partial^2 g+(\partial g)^2.
\end{equation}
Dengan demikian, tensor Riemann memuat informasi mengenai perubahan koneksi dari satu titik ke titik lain dan tidak hanya nilai koneksi pada satu sistem koordinat tertentu.

Tensor Ricci diperoleh melalui kontraksi tensor Riemann,
\begin{equation}
    R_{\mu\nu}=R^\rho_{\ \mu\rho\nu},
    \label{eq:ricci}
\end{equation}
sedangkan skalar kelengkungan diperoleh melalui kontraksi lebih lanjut menggunakan metrik invers,
\begin{equation}
    R=g^{\mu\nu}R_{\mu\nu}.
    \label{eq:ricci-scalar}
\end{equation}
Dari kedua besaran tersebut dibentuk tensor Einstein,
\begin{equation}
    G_{\mu\nu}=R_{\mu\nu}-\frac{1}{2}Rg_{\mu\nu}.
    \label{eq:einstein-tensor}
\end{equation}
Tensor Einstein memenuhi identitas Bianchi terkontraksi,
\begin{equation}
    \nabla_\mu G^{\mu\nu}=0,
    \label{eq:contracted-bianchi}
\end{equation}
yang konsisten dengan konservasi lokal tensor energi-momentum,
\begin{equation}
    \nabla_\mu T^{\mu\nu}=0.
    \label{eq:stress-energy-conservation}
\end{equation}

Hubungan antara geometri ruang-waktu dan distribusi massa-energi dinyatakan oleh persamaan medan Einstein \parencite{Einstein1916GR}:
\begin{equation}
    G_{\mu\nu}
    =\frac{8\pi G}{c^4}T_{\mu\nu},
    \label{eq:efe}
\end{equation}
atau secara eksplisit,
\begin{equation}
    R_{\mu\nu}-\frac{1}{2}Rg_{\mu\nu}
    =\frac{8\pi G}{c^4}T_{\mu\nu}.
\end{equation}
Karena $R_{\mu\nu}$ dan $R$ dibangun dari $g_{\mu\nu}$ beserta turunannya, persamaan medan Einstein merupakan sistem persamaan diferensial nonlinier untuk tensor metrik,
\begin{equation}
    G_{\mu\nu}=G_{\mu\nu}\!\left(g,\partial g,\partial^2 g\right).
\end{equation}

Pada daerah vakum berlaku $T_{\mu\nu}=0$, sehingga
\begin{equation}
    G_{\mu\nu}=0.
\end{equation}
Mengambil trace terhadap persamaan tersebut memberikan
\begin{equation}
    g^{\mu\nu}G_{\mu\nu}
    =R-\frac{1}{2}R\,g^{\mu\nu}g_{\mu\nu}=0.
\end{equation}
Dalam empat dimensi $g^{\mu\nu}g_{\mu\nu}=4$, sehingga
\begin{equation}
    R-2R=0
    \quad\Longrightarrow\quad
    R=0.
\end{equation}
Substitusi kembali ke $G_{\mu\nu}=0$ menghasilkan persamaan medan vakum
\begin{equation}
    R_{\mu\nu}=0.
    \label{eq:vacuum-efe}
\end{equation}
Persamaan~\eqref{eq:vacuum-efe} menjadi titik awal untuk mencari metrik ruang-waktu di luar sumber materi. Dengan asumsi statis dan simetri sferis diperoleh solusi Schwarzschild, sedangkan asumsi stasioner dan simetri aksial untuk sumber berotasi menghasilkan solusi Kerr \parencite{Schwarzschild1916Field,Kerr1963KerrMetric}. Walaupun $R_{\mu\nu}=0$ pada daerah vakum, tensor Riemann secara umum tetap dapat tidak nol, sehingga ruang-waktu vakum tidak harus datar.


% ============================================================
\section{Geodesik dan Transportasi Paralel Spin}
\label{sec:geodesik-spin}

Lintasan partikel uji bermassa dalam ruang-waktu melengkung dinyatakan oleh worldline $x^\mu(\tau)$ dengan empat-kecepatan
\begin{equation}
    u^\mu\equiv\frac{dx^\mu}{d\tau},
    \label{eq:four-velocity}
\end{equation}
di mana $\tau$ adalah waktu proper. Untuk gerak bebas tanpa gaya nongravitasi, lintasan tersebut memenuhi persamaan geodesik \parencite{Schutz2009FirstCourse,Carroll2004SpacetimeGeometry}:
\begin{equation}
    \frac{d^2x^\mu}{d\tau^2}
    +\Gamma^\mu_{\alpha\beta}
    \frac{dx^\alpha}{d\tau}
    \frac{dx^\beta}{d\tau}=0.
    \label{eq:geodesic}
\end{equation}
Dengan menggunakan empat-kecepatan, persamaan yang sama dapat ditulis sebagai
\begin{equation}
    \frac{du^\mu}{d\tau}
    +\Gamma^\mu_{\alpha\beta}u^\alpha u^\beta=0,
\end{equation}
atau secara kovarian,
\begin{equation}
    u^\nu\nabla_\nu u^\mu=0.
    \label{eq:geodesic-covariant}
\end{equation}
Untuk partikel bermassa dan konvensi tanda $(-,+,+,+)$, hubungan $ds^2=-c^2d\tau^2$ memberikan normalisasi
\begin{equation}
    g_{\mu\nu}u^\mu u^\nu=-c^2.
    \label{eq:normalisasi}
\end{equation}

Persamaan geodesik juga dapat diperoleh dari Lagrangian
\begin{equation}
    \mathcal{L}=\frac{1}{2}g_{\mu\nu}\dot{x}^\mu\dot{x}^\nu,
    \label{eq:lagrangian}
\end{equation}
dengan $\dot{x}^\mu\equiv dx^\mu/d\tau$, melalui persamaan Euler--Lagrange
\begin{equation}
    \frac{d}{d\tau}\left(\frac{\partial\mathcal{L}}{\partial\dot{x}^\mu}\right)
    -\frac{\partial\mathcal{L}}{\partial x^\mu}=0.
    \label{eq:euler-lagrange}
\end{equation}
Bentuk Lagrangian ini berguna pada metrik yang memiliki koordinat siklik karena momentum konjugat yang bersesuaian menjadi konstanta gerak.

Orientasi giroskop ideal yang jatuh bebas dinyatakan oleh vektor spin $S^\mu$. Sepanjang worldline geodesik, vektor spin diangkut secara paralel sehingga \parencite{Misner1973Gravitation,Carroll2004SpacetimeGeometry}
\begin{equation}
    \frac{DS^\mu}{D\tau}=0.
    \label{eq:spin-parallel}
\end{equation}
Ekspansi turunan kovarian sepanjang worldline menghasilkan
\begin{equation}
    \frac{DS^\mu}{D\tau}
    =\frac{dS^\mu}{d\tau}
    +\Gamma^\mu_{\nu\rho}S^\nu u^\rho=0,
    \label{eq:spin-transport}
\end{equation}
sehingga evolusi komponen spin dapat dituliskan sebagai sistem ODE orde satu,
\begin{equation}
    \frac{dS^\mu}{d\tau}
    =-\Gamma^\mu_{\nu\rho}S^\nu u^\rho.
    \label{eq:spin-ode}
\end{equation}
Persamaan~\eqref{eq:spin-ode} menunjukkan bahwa penyelesaian evolusi spin memerlukan empat-kecepatan $u^\rho(\tau)$ yang terlebih dahulu diperoleh dari penyelesaian geodesik.

Untuk spin giroskop berlaku kondisi ortogonalitas terhadap empat-kecepatan,
\begin{equation}
    g_{\mu\nu}S^\mu u^\nu=0,
    \label{eq:spin-constraint}
\end{equation}
dan norma spin dipertahankan sepanjang transportasi paralel,
\begin{equation}
    g_{\mu\nu}S^\mu S^\nu=\text{konstan}.
    \label{eq:spin-norm}
\end{equation}
Kedua hubungan tersebut dapat digunakan sebagai pemeriksaan konsistensi dalam integrasi numerik. Secara matematis, rantai ketergantungan yang digunakan dalam penelitian ini dapat diringkas sebagai
\begin{equation}
    g_{\mu\nu}
    \longrightarrow
    \Gamma^\mu_{\alpha\beta}
    \longrightarrow
    x^\mu(\tau),u^\mu(\tau)
    \longrightarrow
    S^\mu(\tau).
\end{equation}
Secara fisis, kondisi $DS^\mu/D\tau=0$ menyatakan bahwa spin tidak mengalami rotasi intrinsik terhadap transportasi paralel. Namun, orientasi spin tetap dapat berubah relatif terhadap kerangka acuan luar; perubahan relatif inilah yang kemudian diamati sebagai presesi geodesik dan presesi Lense--Thirring.


% ============================================================
\subsection{Metrik Schwarzschild}
\label{subsec:metrik-schwarzschild}

Solusi Schwarzschild mendeskripsikan ruang-waktu di sekitar massa statis tak berotasi. Dalam koordinat bola $(t, r, \theta, \phi)$, elemen garisnya diberikan oleh \parencite{Schwarzschild1916Field}:
\begin{equation}
    ds^2 = -\left(1 - \frac{r_s}{r}\right) c^2\,dt^2
         + \left(1 - \frac{r_s}{r}\right)^{-1} dr^2
         + r^2\, d\theta^2
         + r^2 \sin^2\!\theta\, d\phi^2,
    \label{eq:schwarzschild}
\end{equation}
dengan $r_s = 2GM/c^2$ adalah radius Schwarzschild. Metrik ini bersifat diagonal dan menghasilkan efek presesi geodesik, namun tidak mengandung efek \textit{frame-dragging} karena tidak adanya suku campuran $dt\,d\phi$.


\subsection{Metrik Kerr}
\label{subsec:metrik-kerr}

Metrik Kerr merupakan solusi eksak persamaan medan Einstein untuk massa berotasi tanpa muatan listrik, ditemukan oleh \textcite{Kerr1963KerrMetric}. Dalam koordinat Boyer-Lindquist $(t, r, \theta, \phi)$, elemen garisnya dapat dituliskan sebagai:
\begin{equation}
\begin{split}
    ds^2 &= -\left(1 - \frac{r_s r}{\Sigma}\right) c^2\,dt^2
          - \frac{2\,r_s\, r\, a \sin^2\!\theta}{\Sigma}\; c\,dt\,d\phi
          + \frac{\Sigma}{\Delta}\, dr^2
          + \Sigma\, d\theta^2 \\
         &\quad + \left(r^2 + a^2 + \frac{r_s\, r\, a^2 \sin^2\!\theta}{\Sigma}\right)
           \sin^2\!\theta\, d\phi^2,
\end{split}
    \label{eq:kerr-line}
\end{equation}
dengan fungsi-fungsi bantu:
\begin{equation}
    \Sigma \equiv r^2 + a^2 \cos^2\!\theta, \qquad
    \Delta \equiv r^2 - r_s\, r + a^2,
    \label{eq:sigma-delta}
\end{equation}
di mana $r_s = 2GM/c^2$ dan $a = J/(Mc)$ adalah parameter rotasi spesifik, dengan $J$ menyatakan momentum sudut total massa pusat \parencite{Kerr1963KerrMetric,Carroll2004SpacetimeGeometry}.

Dengan menggunakan konvensi koordinat $x^\mu = (x^0, x^1, x^2, x^3) = (ct, r, \theta, \phi)$, tensor metrik Kerr $g_{\mu\nu}$ dapat dituliskan dalam bentuk matriks $4 \times 4$ sebagai:
\begin{equation}
    g_{\mu\nu} =
    \begin{pmatrix}
        -\left(1 - \dfrac{r_s r}{\Sigma}\right) & 0 & 0 & -\dfrac{r_s\, r\, a \sin^2\!\theta}{\Sigma} \\[12pt]
        0 & \dfrac{\Sigma}{\Delta} & 0 & 0 \\[12pt]
        0 & 0 & \Sigma & 0 \\[12pt]
        -\dfrac{r_s\, r\, a \sin^2\!\theta}{\Sigma} & 0 & 0
        & \left(r^2 + a^2 + \dfrac{r_s\, r\, a^2 \sin^2\!\theta}{\Sigma}\right)\sin^2\!\theta
    \end{pmatrix}.
    \label{eq:kerr-matrix}
\end{equation}

Beberapa sifat penting dari matriks metrik ini:
\begin{enumerate}[label=(\roman*)]
    \item Komponen $g_{t\phi} = g_{\phi t} = -r_s\, r\, a \sin^2\!\theta\,/\,\Sigma$ tidak bernilai nol jika $a \neq 0$. Suku campuran inilah yang bertanggung jawab atas efek \textit{frame-dragging}.
    \item Pada limit $a \to 0$, matriks di atas tereduksi menjadi metrik Schwarzschild diagonal (Persamaan~\ref{eq:schwarzschild}).
    \item Komponen $g_{rr}$ memiliki singularitas koordinat pada $\Delta = 0$, yang terjadi pada horizon peristiwa.
\end{enumerate}

Metrik invers $g^{\mu\nu}$ diperlukan untuk menghitung simbol Christoffel. Dengan memanfaatkan struktur blok-diagonal ($2 \times 2$ pada sub-ruang $t$-$\phi$ dan diagonal pada $r$-$\theta$), komponen-komponen metrik invers diperoleh sebagai:
\begin{equation}
\begin{aligned}
    g^{tt} &= -\frac{1}{\Delta}\left(r^2 + a^2 + \frac{r_s\, r\, a^2 \sin^2\!\theta}{\Sigma}\right), &\qquad
    g^{t\phi} &= -\frac{r_s\, r\, a}{\Sigma\,\Delta}, \\[6pt]
    g^{rr} &= \frac{\Delta}{\Sigma}, &\qquad
    g^{\theta\theta} &= \frac{1}{\Sigma}, \\[6pt]
    g^{\phi\phi} &= \frac{1}{\Delta\sin^2\!\theta}\left(1 - \frac{r_s\, r}{\Sigma}\right), &\qquad
    & \text{komponen lain} = 0.
\end{aligned}
    \label{eq:kerr-inverse}
\end{equation}


\subsection{Contoh Perhitungan Simbol Christoffel Metrik Kerr}
\label{subsec:christoffel-kerr-contoh}

Perhitungan simbol Christoffel untuk metrik Kerr menghasilkan sejumlah besar komponen yang tidak bernilai nol. Pada bagian ini disajikan derivasi eksplisit untuk tiga komponen representatif yang relevan dengan efek \textit{frame-dragging}.

\paragraph{Contoh 1: $\Gamma^t_{tr}$.}
Menggunakan Persamaan~\eqref{eq:christoffel-def} dengan $\mu = t$, $\nu = t$, $\rho = r$:
\begin{equation}
    \Gamma^t_{tr} = \frac{1}{2}\,g^{t\sigma}
    \left(\partial_r g_{\sigma t} + \partial_t g_{\sigma r} - \partial_\sigma g_{tr}\right).
\end{equation}
Karena metrik tidak bergantung pada $t$ (stasioner), suku $\partial_t g_{\sigma r} = 0$. Kontribusi non-nol berasal dari $\sigma = t$ dan $\sigma = \phi$:
\begin{equation}
    \Gamma^t_{tr} = \frac{1}{2}\,g^{tt}\,\partial_r g_{tt}
                  + \frac{1}{2}\,g^{t\phi}\,\partial_r g_{\phi t}.
    \label{eq:christoffel-ttr}
\end{equation}
Turunan parsial $\partial_r g_{tt}$ dan $\partial_r g_{t\phi}$ dapat dihitung dari Persamaan~\eqref{eq:kerr-matrix}. Sebagai contoh:
\begin{equation}
    \frac{\partial g_{tt}}{\partial r}
    = \frac{\partial}{\partial r}\left[-\left(1 - \frac{r_s\, r}{\Sigma}\right)\right]
    = -\frac{r_s(r^2 - a^2\cos^2\!\theta)}{\Sigma^2},
    \label{eq:dgtt-dr}
\end{equation}
mengingat $\partial_r \Sigma = 2r$. Substitusi ke Persamaan~\eqref{eq:christoffel-ttr} bersama dengan ekspresi $g^{tt}$ dan $g^{t\phi}$ dari Persamaan~\eqref{eq:kerr-inverse} menghasilkan:
\begin{equation}
    \Gamma^t_{tr} = \frac{r_s(r^2 - a^2\cos^2\!\theta)(r^2 + a^2)}{2\,\Sigma^2\,\Delta}
    + \frac{r_s^2\, r\, a^2\sin^2\!\theta\,(r^2 - a^2\cos^2\!\theta)}{2\,\Sigma^3\,\Delta}.
    \label{eq:gamma-ttr-result}
\end{equation}

\paragraph{Contoh 2: $\Gamma^\phi_{tr}$.}
Komponen ini menggambarkan hubungan antara gerak radial dan gerak azimut akibat rotasi sumber:
\begin{equation}
    \Gamma^\phi_{tr} = \frac{1}{2}\,g^{\phi\sigma}
    \left(\partial_r g_{\sigma t} + \partial_t g_{\sigma r} - \partial_\sigma g_{tr}\right).
\end{equation}
Dengan $\partial_t = 0$ dan kontribusi dari $\sigma = t$ serta $\sigma = \phi$:
\begin{equation}
    \Gamma^\phi_{tr} = \frac{1}{2}\,g^{\phi t}\,\partial_r g_{tt}
                     + \frac{1}{2}\,g^{\phi\phi}\,\partial_r g_{\phi t}.
    \label{eq:christoffel-phitr}
\end{equation}
Perhatikan bahwa $g^{\phi t} = g^{t\phi}$. Substitusi komponen metrik invers dan turunan yang sesuai memberikan ekspresi dalam $r$, $\theta$, $\Sigma$, dan $\Delta$. Komponen ini bernilai nol pada limit $a \to 0$, yang menegaskan bahwa hubungan $r$--$\phi$ tersebut berasal dari rotasi massa.

\paragraph{Contoh 3: $\Gamma^t_{\theta\phi}$.}
Komponen ini menunjukkan hubungan antara gerak polar dan azimut melalui koordinat waktu:
\begin{equation}
    \Gamma^t_{\theta\phi} = \frac{1}{2}\,g^{t\sigma}
    \left(\partial_\phi g_{\sigma\theta} + \partial_\theta g_{\sigma\phi} - \partial_\sigma g_{\theta\phi}\right).
\end{equation}
Karena metrik tidak bergantung pada $\phi$ dan $g_{r\theta} = g_{r\phi} = 0$, maka
\begin{equation}
    \Gamma^t_{\theta\phi} = \frac{1}{2}\,g^{tt}\,\partial_\theta g_{t\phi}
                          + \frac{1}{2}\,g^{t\phi}\,\partial_\theta g_{\phi\phi}.
    \label{eq:christoffel-tthetaphi}
\end{equation}
Turunan $\partial_\theta g_{t\phi}$ diberikan oleh
\begin{equation}
    \frac{\partial g_{t\phi}}{\partial\theta}
    = -\frac{r_s\, r\, a}{\Sigma^2}
    \left(2\Sigma\sin\theta\cos\theta + 2a^2\sin^3\!\theta\cos\theta\right),
\end{equation}
yang menunjukkan ketergantungan eksplisit pada sudut polar $\theta$. Komponen ini juga lenyap pada limit $a \to 0$.

Perhitungan lengkap seluruh komponen Christoffel metrik Kerr menghasilkan puluhan suku yang kompleks. Dalam praktik, perhitungan ini dapat diverifikasi dan dilakukan secara efisien menggunakan pustaka komputasi simbolik seperti \texttt{EinsteinPy} \parencite{Bapat2020EinsteinPy}.


% ============================================================
\section{Efek Relativistik di Sekitar Bumi}
\label{sec:efek-relativistik}

\subsection{Efek Geodesik (Presesi de Sitter)}
\label{subsec:efek-geodesik}

Efek geodesik adalah presesi orientasi giroskop yang bergerak dalam medan gravitasi akibat kelengkungan ruang-waktu, bahkan tanpa rotasi sumber. Efek ini muncul dari transportasi paralel vektor spin sepanjang geodesik pada metrik Schwarzschild. Dalam aproksimasi medan lemah dan orbit sirkular, laju presesi geodesik diberikan oleh \parencite{Will2014Confrontation}:
\begin{equation}
    \bm{\Omega}_{\text{geodetik}} = \frac{3}{2}\,\frac{GM}{c^2\,r^3}\; \mathbf{r} \times \mathbf{v},
    \label{eq:omega-geodetic}
\end{equation}
dengan $\mathbf{r}$ adalah vektor posisi dan $\mathbf{v}$ adalah kecepatan orbit. Untuk orbit sirkular dengan $v = \sqrt{GM/r}$, magnitudo laju presesi menjadi:
\begin{equation}
    |\bm{\Omega}_{\text{geodetik}}| = \frac{3}{2}\,\frac{(GM)^{3/2}}{c^2\, r^{5/2}}.
    \label{eq:omega-geodetic-circular}
\end{equation}


\subsection{Efek Frame-Dragging (Presesi Lense-Thirring)}
\label{subsec:frame-dragging}

Efek Lense-Thirring muncul karena rotasi massa sumber yang menyeret ruang-waktu di sekitarnya. Secara matematis, efek ini berasal dari komponen campuran $g_{t\phi} \neq 0$ pada metrik Kerr yang tidak ada pada metrik Schwarzschild. Bentuk vektor umum laju presesi Lense-Thirring pada jarak $\mathbf{r}$ dari sumber bermomen sudut $\mathbf{J}$ diberikan oleh \parencite{Lense1918LenseThirring}:
\begin{equation}
    \bm{\Omega}_{\text{LT}}(\mathbf{r})
    = \frac{G}{c^2\, r^3}\left[3\hat{\mathbf{r}}\,(\hat{\mathbf{r}} \cdot \mathbf{J}) - \mathbf{J}\right],
    \label{eq:omega-lt-general}
\end{equation}
dengan $\hat{\mathbf{r}} = \mathbf{r}/r$. Untuk orbit polar (relevan dengan konfigurasi GP-B), ekspresi ini dapat disederhanakan menjadi:
\begin{equation}
    |\bm{\Omega}_{\text{LT}}| \approx \frac{2G\,|\mathbf{J}|}{c^2\, r^3}.
    \label{eq:omega-lt-approx}
\end{equation}
Derivasi lengkap ekspresi ini dari metrik Kerr melalui reduksi medan lemah disajikan pada Bab~\ref{chap:metode}.

Laju presesi total yang dialami giroskop merupakan superposisi kedua kontribusi:
\begin{equation}
    \bm{\Omega} = \bm{\Omega}_{\text{geodetik}} + \bm{\Omega}_{\text{LT}}.
    \label{eq:omega-total}
\end{equation}


% ============================================================
\section{Eksperimen Gravity Probe B}
\label{sec:gravity-probe-b}

Eksperimen \textit{Gravity Probe B} (GP-B) merupakan misi satelit yang dirancang oleh Stanford University dan NASA untuk menguji dua prediksi Relativitas Umum: efek geodesik dan efek \textit{frame-dragging} \parencite{Turneaure1986GravityProbeBApproach}. Misi ini diluncurkan pada 20 April 2004 dan mengumpulkan data selama lebih dari satu tahun.

Instrumen utama GP-B terdiri dari empat giroskop kuarsa superpresisi yang mengorbit Bumi dalam lintasan polar pada ketinggian $\sim$642~km (radius orbit $r \approx 7{,}013 \times 10^6$~m). Orientasi referensi diarahkan ke bintang panduan IM Pegasi. Parameter fisik yang relevan ditampilkan pada Tabel~\ref{tab:param-gpb}.

\begin{table}[ht]
\centering
\caption{Parameter fisik yang digunakan dalam analisis GP-B.}
\label{tab:param-gpb}
\begin{tabular}{lll}
\toprule
\textbf{Parameter} & \textbf{Simbol} & \textbf{Nilai} \\
\midrule
Massa Bumi             & $M$    & $5{,}972 \times 10^{24}$~kg \\
Momentum sudut Bumi    & $J$    & $5{,}86 \times 10^{33}$~kg\,m$^2$\,s$^{-1}$ \\
Radius orbit           & $r$    & $7{,}013 \times 10^{6}$~m \\
Radius Schwarzschild   & $r_s$  & $8{,}87 \times 10^{-3}$~m \\
Parameter rotasi       & $a$    & $3{,}98 \times 10^{-3}$~m \\
Konstanta gravitasi    & $G$    & $6{,}674 \times 10^{-11}$~m$^3$\,kg$^{-1}$\,s$^{-2}$ \\
Kecepatan cahaya       & $c$    & $2{,}998 \times 10^{8}$~m\,s$^{-1}$ \\
\bottomrule
\end{tabular}
\end{table}

Hasil akhir pengukuran GP-B yang dipublikasikan oleh \textcite{Everitt2011GravityProbeB} adalah:
\begin{equation}
    \Omega_{\text{geodetik}}^{\text{GP-B}} = 6602 \pm 18~\text{mas/yr}, \qquad
    \Omega_{\text{LT}}^{\text{GP-B}} = 37{,}2 \pm 7{,}2~\text{mas/yr},
    \label{eq:gpb-results}
\end{equation}
di mana 1~mas (milliarcsecond) $= 4{,}848 \times 10^{-9}$~rad. Kedua hasil ini konsisten dengan prediksi Relativitas Umum. Konversi satuan dari rad/s ke mas/yr menggunakan faktor:
\begin{equation}
    1~\text{rad/s} = 206\,264\,806{,}247 \times 3{,}15576 \times 10^7~\text{mas/yr}
    \approx 6{,}51 \times 10^{15}~\text{mas/yr}.
    \label{eq:konversi-masyr}
\end{equation}

Selain GP-B, efek Lense-Thirring juga dikonfirmasi oleh analisis orbit satelit LAGEOS dengan ketidakpastian sekitar 10\% \parencite{CiufoliniPavlis2004_LAGEOS}. Berbagai pendekatan alternatif, termasuk model gravitasi vektor \parencite{BeheraBarik2020} dan teori modifikasi paritas \parencite{Mao2006,Nakamura2018}, telah diajukan namun belum menunjukkan deviasi signifikan dari GR pada tingkat presisi eksperimen saat ini.


% ============================================================
\section{Pustaka Komputasi Python}
\label{sec:pustaka-python}

Penyelesaian numerik persamaan geodesik dan transportasi paralel spin memerlukan implementasi komputasi yang efisien. Penelitian ini memanfaatkan pustaka Python berikut sebagai alat bantu:
\begin{enumerate}[label=(\roman*)]
    \item \textbf{EinsteinPy} \parencite{Bapat2020EinsteinPy}: pustaka sumber terbuka untuk komputasi tensorial GR, menyediakan kelas untuk mendefinisikan metrik, menghitung simbol Christoffel, dan mengintegrasikan geodesik.
    \item \textbf{FANTASY} \parencite{Christian2020FANTASY}: integrator geodesik simplektik dengan diferensiasi otomatis, menawarkan stabilitas numerik jangka panjang.
    \item \textbf{KerrGeoPy} \parencite{Park2024KerrGeoPy}: pustaka khusus untuk geodesik pada metrik Kerr.
    \item \textbf{SciPy} \parencite{SciPy2020}: pustaka komputasi ilmiah yang menyediakan solver ODE (\texttt{solve\_ivp}) dengan berbagai metode termasuk Runge-Kutta orde 4/5 (RK45).
    \item \textbf{Matplotlib} \parencite{Matplotlib2007}: pustaka visualisasi data untuk menghasilkan grafik evolusi spin dan analisis kestabilan numerik.
\end{enumerate}

Pustaka-pustaka ini digunakan untuk memverifikasi hasil derivasi analitik dan menghasilkan solusi numerik yang disajikan pada Bab~\ref{chap:hasil}. Penjelasan detail mengenai penggunaan metode numerik disajikan pada Bab~\ref{chap:metode}.